\documentclass[10pt,a4paper]{amsart}
\usepackage[margin=19mm]{geometry}
\usepackage{amsmath,amssymb,mathtools}
\usepackage[scr=rsfs,cal=euler]{mathalfa}
\usepackage{xfrac}
\usepackage[unicode,hidelinks,bookmarks=false]{hyperref}
\newcommand{\ie}{i.e.\ }
\newcommand{\cf}{cf.\ }
\newcommand{\resp}{respectively\ }
\newcommand{\Subsection}{\subsection}
\providecommand{\nobreakdash}{\nobreak\hbox{-}}
\setlength{\emergencystretch}{2em}
\multlinegap0pt
\usepackage{tikz}
\usetikzlibrary{cd}
\usepackage{amssymb}
\usepackage{xcolor}
\usepackage{braket}
\usepackage[shortlabels]{enumitem}
\usepackage{tikz}
\newtheorem{lemma}{Lemma}
\newtheorem{proposition}{Proposition}
\newcommand{\CC}{\mathbb{C}}
\newcommand{\Hrel}{\mathsf{H}_{\mathrm{rel}}}
\newcommand{\VV}{\mathbb{V}}
\newcommand{\be}{\boldsymbol{e}}
\newcommand\ext[1]{\scalebox{.85}{$\bigwedge^{\!#1}$}}
\newcommand\sym[1]{\scalebox{.85}{$\bigodot^{\!#1}$}}
\title{The spectrum of the bosonic ambitwistor string revisited}
\author{Jos\'e M. Figueroa-O'Farrill, Girish S. Vishwa}
\date{Source: arXiv:2606.05500v2 (CC BY 4.0)}
\begin{document}
\maketitle

\noindent\emph{Source and licence.} The spectrum of the bosonic ambitwistor string revisited. Version arXiv:2606.05500v2 (CC BY 4.0), distributed by arXiv under the Creative Commons Attribution 4.0 International licence, http://creativecommons.org/licenses/by/4.0/, as recorded in the arXiv licence metadata for this exact version and shown on the abstract page https://arxiv.org/abs/2606.05500v2. Full licence notice: \texttt{licenses/arxiv-2606.05500-CC-BY-4.0.txt}.\par\smallskip

\noindent\textit{Editorial scope.} Counted unit: the article's proposition prop:poincare-duality (source0.tex lines 2991--2994). The article's complete original proof of it is delivered with it (source0.tex lines 2996--3068, one \texttt{proof} environment of 73 lines). No mathematics is written, repaired or re-derived: the statement and the proof are the article's own lines, in the article's own notation, and no retained line is substituted or re-flowed.  Retained uncounted as the source-local context the proof needs: lines 1190--1203; lines 2808--2816.  Nothing was omitted from any retained span, so the package carries no omission record.  No retained line contains a citation, so the wrapper carries no bibliography and no bibliography entry.  The retained lines contain the article's own diagram code, which the wrapper typesets with the standard tikz packages; no image file is delivered and the package declares no asset dependency.  Editorial formatting only, in addition: an AMS \texttt{amsart} preamble that restates the article's own declarations and macros used by the retained lines, namely the environments \texttt{lemma}, \texttt{proposition} on separate counters, together with the standard packages those lines need.  The retained environments print in this document as Lemma 1, Proposition 1 in order of appearance, whereas the article prints the same items with the numbering of its own full document.  The complete native source of the article as distributed in the arXiv e-print for this exact version is bundled unchanged as \texttt{sources/202121/source0.tex}.
\begin{lemma}
  \label{lem:iso-h2rel-skewsymmetric}
  The canonical surjection $\pi : p^\circ \to V^\top$ induces an
  isomorphism of $H$-modules
  \begin{equation}
    \dfrac{\ext{2}p^\circ}{\ell_p \wedge p^\circ} \cong \ext{2}V^\top
  \end{equation}
  and hence a short exact sequence
  \begin{equation}
    \begin{tikzcd}
      0 \arrow[r] & \ell_p \wedge p^\circ \arrow[r] & \ext{2}p^\circ \arrow[r] & \ext{2} V^\top \arrow[r] & 0.
    \end{tikzcd}
  \end{equation}
\end{lemma}

\begin{equation}
  \label{eq:gram-det}
  \bar\eta(\bar u_1 \wedge \bar v_1, \bar u_2 \wedge \bar v_2) =
  \det
  \begin{pmatrix}
    \eta(u_1,u_2) & \eta(u_1,v_2) \\
    \eta(v_1,u_2) & \eta(v_1,v_2)
  \end{pmatrix}
\end{equation}

\begin{proposition}[Poincaré duality]
  \label{prop:poincare-duality}
  There is an $H$-module isomorphism $\Hrel^2(p) \cong \Hrel^3(p)^*$.
\end{proposition}

\begin{proof}
  This is equivalent to the existence of a non-degenerate
  $H$-equivariant pairing
  \begin{equation}
    \Hrel^2(p) \times \Hrel^3(p) \to \CC
  \end{equation}
  which we will now exhibit.  We have that
  \begin{equation}
    \Hrel^2(p) \cong \dfrac{\ext{2} p^\perp}{L_p \wedge p^\perp} \oplus \dfrac{\sym{2} p^\perp}{L_p \odot L_p}
  \end{equation}
  whereas
  \begin{equation}
    \Hrel^3(p) \cong \dfrac{\ext{2} p^\perp}{L_p \wedge p^\perp} \oplus \dfrac{p^\perp \odot \VV}{L_p \odot \VV}.
  \end{equation}
  It follows from Lemma~\ref{lem:iso-h2rel-skewsymmetric} that
  \begin{equation}
    \dfrac{\ext{2}p^\perp}{L_p \wedge p^\perp} \cong \ext{2} \VV^\top,
  \end{equation}
  which inherits a positive-definite $H$-invariant inner product from
  the euclidean inner product on $\VV^\top$ via the Gram
  determinant~\eqref{eq:gram-det}.  So we must now exhibit a
  non-degenerate pairing between
  \begin{equation}
     \dfrac{\sym{2} p^\perp}{L_p \odot L_p} \qquad\text{and}\qquad
     \dfrac{p^\perp \odot \VV}{L_p \odot \VV}.
  \end{equation}
  The complex-bilinear inner product on $\sym{2}\VV$ induced by $\eta$
  via the symmetric analogue of the Gram determinant gives a dual
  pairing
  \begin{equation}
    \sym{2}\VV \times \sym{2}\VV \to \CC,
  \end{equation}
  which we may restrict to a bilinear map
  \begin{equation}
    \sym{2} p^\perp \times p^\perp \odot \VV \to \CC.
  \end{equation}
  Explicitly, if $w_1,w_2,w_3 \in p^\perp$ and $v \in \VV$, this map
  sends
  \begin{equation}
    (w_1 \odot w_2, w_3 \odot v) \mapsto \eta(w_1,w_3) \eta(w_2,v) + \eta(w_2,w_3) \eta(w_1,v).
  \end{equation}
  Clearly if $w_1, w_2 \in L_p$, then this is zero for all $w_3 \in
  p^\perp$ and $v \in \VV$ and, similarly, if $w_3 \in L_p$, then this
  is zero for all $w_1,w_2 \in p^\perp$ and any $v \in \VV$.  Therefore
  we see that the above bilinear map descends to a bilinear map
  \begin{equation}
    \left<-,-\right> : \dfrac{\sym{2} p^\perp}{L_p \odot L_p} \times \dfrac{p^\perp \odot \VV}{L_p \odot \VV} \to \CC.
  \end{equation}
  We can show that this pairing is non-degenerate by explicitly
  writing it down in a basis.  A Witt frame $\be_+,\be_i,
  \be_-$ for $\VV$ gives a natural basis $\be_+ \odot \be_+, \be_+
  \odot \be_i$, $\be_i \odot \be_j$ (for $i\leq j$) for $\sym{2}p^\perp$
  and $\be_+\odot \be_+, \be_+ \odot \be_i, \be_+\odot \be_-, \be_i
  \odot \be_j, \be_i \odot \be_-$ (for $i\leq j$) for $p^\perp \odot
  \VV$. Letting bars denote the projection to the relevant quotient, a
  basis for $\sym{2}p^\perp/L_p \odot L_p$ is given by
  $\overline{\be_+ \odot \be_i}, \overline{\be_i \odot \be_j}$ (for 
  $i\leq j$), whereas a basis for $(p^\perp \odot \VV)/(L_p \odot
  \VV)$ is given by $\overline{\be_- \odot \be_i}, \overline{\be_i
    \odot \be_j}$ (for $i\leq j$).  The above pairing on these basis
  elements is given by
  \begin{equation}
    \begin{split}
      \left<\overline{\be_+ \odot \be_i}, \overline{\be_- \odot
          \be_j}\right> &= \eta(\be_+,\be_-) \eta(\be_i, \be_j) = \delta_{ij}\\
      \left<\overline{\be_i \odot \be_j}, \overline{\be_k \odot
          \be_\ell}\right> &= \eta(\be_i, \be_k) \eta(\be_j,\be_\ell)
      +  \eta(\be_i, \be_\ell) \eta(\be_j,\be_k) = \delta_{ik}\delta_{j\ell} +  \delta_{i\ell} \delta_{jk},
  \end{split}
\end{equation}
which is clearly non-degenerate.  We conclude that $\Hrel^2(p)$ and
$\Hrel^3(p)$ are dual $H$-modules.
\end{proof}

\end{document}
